\documentclass[12pt]{article}
\usepackage[margin=1in]{geometry}
\usepackage{amsmath,amssymb,amsthm,booktabs,natbib,microtype}
\newtheorem{theorem}{Theorem}
\newtheorem{proposition}[theorem]{Proposition}
\newtheorem{corollary}[theorem]{Corollary}
\newtheorem{remark}[theorem]{Remark}
\newcommand{\norm}[1]{\lVert#1\rVert}
\title{Common-Coordinate Transport of Riesz Clouds\\\large Summable Resolvent Defects and Fixed-Rank Limits}
\author{Prime Dynamics Theory Program}
\date{July 2026}

\begin{document}
\maketitle
\begin{abstract}
RH-162--168 give fixed-scale packet-to-Riesz certificates.  An all-level
interface also needs a coherent comparison of Riesz clouds living at
different mesh and noise scales.  After embedding two consecutive operators
in a common Hilbert space, we prove
\[
 \norm{\Pi_{j+1}-\Pi_j}
 \le\frac{|\Gamma|}{2\pi}M_{j+1}M_j
 \norm{A_{j+1}-A_j},
\]
where one contour $\Gamma$ lies in both resolvent sets and $M_j$ are its
resolvent bounds.  A step bound below one yields an invertible range
transport with explicit inverse norm.  If the step bounds are summable, the
Riesz projections converge in operator norm to an idempotent of the same
finite rank.  A 512-pair Hermitian audit has no inequality failure.  The
theorem is restricted to fixed rank; a moving determinant cloud has growing
rank and requires the shellwise correction developed next.
\end{abstract}

\section{The common-coordinate input}
Operators at different Ulam dimensions cannot be subtracted until maps into
one ambient space have been chosen.  Let $\mathcal H_j$ be the scale spaces,
let $J_j:\mathcal H_j\to\mathcal K$ be isometries, and suppose extensions
$\widetilde A_j\in\mathcal B(\mathcal K)$ have been specified.  Their action
outside $J_j\mathcal H_j$ is part of the construction; arbitrary zero
extension can introduce spurious spectrum on a contour.

For the theorem, write $A_j$ for these common-space operators.  Let one
rectifiable Jordan curve $\Gamma$ satisfy
\[
 \Gamma\subset\rho(A_j)\cap\rho(A_{j+1}),\qquad
 \sup_\Gamma\norm{(z-A_j)^{-1}}\le M_j.
\]
Define
\[
 \Pi_j=\frac1{2\pi i}\int_\Gamma(z-A_j)^{-1}\,dz.
\]

\section{Consecutive transport theorem}
\begin{theorem}[Common-contour Riesz transport]\label{thm:step}
If $\norm{A_{j+1}-A_j}\le\eta_j$, then
\[
 \boxed{
 \norm{\Pi_{j+1}-\Pi_j}
 \le \beta_j:=\frac{|\Gamma|}{2\pi}M_{j+1}M_j\eta_j.
 \ }
\]
\end{theorem}
\begin{proof}
The resolvent identity gives
\[
 (z-A_{j+1})^{-1}-(z-A_j)^{-1}
 =(z-A_{j+1})^{-1}(A_{j+1}-A_j)(z-A_j)^{-1}.
\]
Integrate around $\Gamma$ and take norms \citep{Kato1995}.
\end{proof}

\begin{corollary}[Stable range map]\label{cor:range}
Suppose $\Pi_j$ and $\Pi_{j+1}$ have equal finite rank and $\beta_j<1$.
Then
\[
 \Pi_{j+1}|_{\operatorname{ran}\Pi_j}:
 \operatorname{ran}\Pi_j\longrightarrow\operatorname{ran}\Pi_{j+1}
\]
is invertible, and its inverse norm is at most $(1-\beta_j)^{-1}$.
\end{corollary}
\begin{proof}
For $x\in\operatorname{ran}\Pi_j$,
\[
 \norm{\Pi_{j+1}x}\ge
 \norm x-\norm{(\Pi_{j+1}-\Pi_j)x}
 \ge(1-\beta_j)\norm x.
\]
Thus the map is injective; equal finite dimensions give surjectivity and the
inverse bound.
\end{proof}

\section{Summable fixed-rank limit}
\begin{theorem}[Norm-convergent Riesz chain]\label{thm:sum}
Assume the same finite rank $r$ is enclosed at every scale and
\[
 \sum_{j=1}^{\infty}\beta_j<\infty.
\]
Then $\Pi_j$ converges in operator norm to a bounded idempotent $\Pi_\infty$
of rank $r$.  The tail obeys
\[
 \norm{\Pi_\infty-\Pi_n}\le\sum_{j=n}^{\infty}\beta_j.
\]
\end{theorem}
\begin{proof}
The telescoping inequality makes $(\Pi_j)$ norm-Cauchy.  Completeness of
$\mathcal B(\mathcal K)$ gives a limit, and multiplication is norm
continuous, so $\Pi_\infty^2=\Pi_\infty$.  For sufficiently large $j$,
$\norm{\Pi_j-\Pi_\infty}<1$; the close-projection restriction argument gives
equal rank.  The tail estimate follows from telescoping.
\end{proof}

Summability is stronger than $\eta_j\to0$.  Resolvent conditioning can grow
fast enough that $M_jM_{j+1}\eta_j$ is not summable even when operators look
close entrywise.  This is the all-level analogue of the fixed-scale
pseudospectral wall.

\section{Moving contours and homotopy compatibility}
Consecutive physical shells may use different contours.  They can be
compared by Theorem~\ref{thm:step} only after a common representative is
chosen.

\begin{proposition}[Contour deformation ledger]
Let $\Gamma_0$ and $\Gamma_1$ be homologous closed contours in one connected
open subset of $\rho(A)$, enclosing the same spectral component.  Then their
Riesz integrals are equal.  Consequently, if both consecutive operators are
resolvent-free throughout a common deformation corridor, either boundary may
be used in the transport estimate.
\end{proposition}
\begin{proof}
The operator-valued resolvent is holomorphic on the corridor.  Cauchy's
theorem makes the two contour integrals equal \citep{Kato1995}.
\end{proof}

This adds a concrete schedule obligation.  Merely labeling two independently
chosen clouds by the same shell index does not prove they enclose the same
spectral component; a joint resolvent corridor or an equivalent coefficient
identification is required before telescoping.

\section{Vanishing steps are not enough}
The summability hypothesis in Theorem~\ref{thm:sum} cannot be replaced by
$\beta_j\to0$ in this proof architecture.

\begin{proposition}[Slow rotating packet witness]
There is a sequence of rank-one orthogonal projections on $\mathbb C^2$ such
that
\[
 \norm{P_{j+1}-P_j}\longrightarrow0
\]
but $(P_j)$ does not converge.
\end{proposition}
\begin{proof}
Let
\[
 \theta_j=\sum_{k=1}^{j}\frac1k,qquad
 u_j=(\cos\theta_j,\sin\theta_j),\qquad P_j=u_ju_j^*.
\]
Then
$\norm{P_{j+1}-P_j}=|\sin(1/(j+1))|\to0$.  The angle $\theta_j$ keeps
traversing the projective circle because its increments tend to zero while
their sum diverges.  Hence the ranges visit separated angular sectors
infinitely often and do not converge.
\end{proof}

Taking $A_j=P_j$, a fixed circle around the eigenvalue one gives uniformly
bounded resolvents and Riesz projection $P_j$.  Thus even an exact uniformly
gapped spectral model can drift forever when the common-coordinate defects
are nonsummable.

\section{Conditioning of the composed range transport}
Let
\[
 L_j=\Pi_{j+1}|_{\operatorname{ran}\Pi_j}
\]
be the step maps from Corollary~\ref{cor:range}.  Their composition transports
one initial cloud frame through the chain without making a new gauge choice
at every scale.

\begin{proposition}[Finite cumulative inverse condition]
If $\beta_j\le\beta_*<1$ and $\sum_j\beta_j<\infty$, then
\[
 \prod_j\norm{L_j^{-1}}
 \le\prod_j(1-\beta_j)^{-1}<\infty.
\]
More explicitly,
\[
 \log\prod_{j=n}^{N}(1-\beta_j)^{-1}
 \le\frac1{1-\beta_*}\sum_{j=n}^{N}\beta_j.
\]
\end{proposition}
\begin{proof}
Corollary~\ref{cor:range} gives the first finite-product inequality.  For
$0\le x\le\beta_*$,
$-\log(1-x)\le x/(1-\beta_*)$.  Sum this scalar inequality and exponentiate.
\end{proof}

Thus the same summability that yields a projector limit also prevents the
canonical stepwise coordinates from acquiring infinite inverse condition.
This is stronger information than pairwise rank agreement and is relevant
when shell coefficients or marked vectors must be transported coherently.

\section{Audit and boundary}
The finite audit rotates a three-dimensional spectral packet inside eight-
dimensional Hermitian matrices while keeping a unit-circle spectral gap.  It
evaluates the exact normal resolvent bound, the operator defect, and the
projector distance for 512 pairs.  No distance exceeds
Theorem~\ref{thm:step}'s envelope.

RH-169 proves the correct fixed-rank common-coordinate theorem.  It does not
construct the physical embeddings or extensions, prove summability, or
handle the increasing degree of the RH-80 cloud.  Therefore physical R,
Gate A, and all later Hilbert--Polya interfaces remain open.

\bibliographystyle{plainnat}
\bibliography{references}
\end{document}
